% !TEX program = xelatex
% Version grands caractères (paysage, ~20pt, une page par exercice) de
% "Séquence 2 - Feuille d'exercices - Vecteurs.tex" (Feuille 2), pour une
% élève déficiente visuelle. Document autonome (n'utilise PAS CLT2627.sty),
% sur le même principe que "Séquence 2 - Cours - Vecteurs - Grands
% Caractères.tex" : police sans-serif, texte non justifié, interligne
% large, paysage (comme le cours), et chaque exercice tient sur sa propre
% page pour rester lisible à cette taille. Quand l'exercice a une figure,
% l'énoncé occupe la colonne de gauche et la figure la colonne de droite
% (voir \GCexo plus bas) : la page paysage étant large mais basse, c'est
% ce qui laisse le plus de place à la figure tout en gardant l'ensemble
% sur une seule page.
%
% Ne couvre que les Objectifs 1 à 4 / Exercices 1 à 16 : la feuille source
% a un \end{document} prématuré juste après l'exercice 16 qui rend tout le
% reste (Objectifs 5-7, exercices 18-26) inactif dans le PDF normal ; cette
% version grands caractères reste donc volontairement alignée sur ce qui
% est réellement diffusé aux élèves.
\documentclass[a4paper,landscape]{article}

\usepackage{fontspec}
\usepackage[french]{babel}
\usepackage{geometry}
\geometry{a4paper,landscape,top=1.4cm,bottom=1.6cm,left=1.8cm,right=1.8cm,footskip=1cm}

\setmainfont{Helvetica Neue}

\usepackage{amsmath,amssymb}
\usepackage{xcolor}
\usepackage[most]{tcolorbox}
\tcbuselibrary{skins,breakable}
\usepackage[inline]{enumitem}
\usepackage{ragged2e}
\usepackage{fancyhdr}
\usepackage{lastpage}
\usepackage{graphicx}
\usepackage{adjustbox}
\usepackage{pifont}
% Badges « Exercice n » en cursive Playwrite FR Moderne (police du site),
% chargée à 85 % de la taille courante (XeLaTeX ou LuaLaTeX ; gras sous pdfLaTeX).
\makeatletter
\ifdefined\XeTeXrevision\def\ExoFont@ok{}\fi
\ifdefined\directlua\def\ExoFont@ok{}\fi
\ifdefined\ExoFont@ok
\newcommand{\ExoFont}{\font\ExoFont@f="[PlaywriteFRModerne.ttf]" at \dimexpr\f@size pt*85/100\relax\ExoFont@f}
\else
\newcommand{\ExoFont}{\bfseries}
\fi
\makeatother
% \ExoTexte{...} : texte du badge en cursive, avec la hauteur d'une majuscule de la
% police courante (la cursive est plus haute et ferait grandir chaque badge).
\newcommand{\ExoTexte}[1]{\raisebox{0pt}[\fontcharht\font`E][0pt]{\ExoFont #1}}
\usepackage{frcursive}
% Correctif du bug de crénage de frcursive sous XeLaTeX (repris de
% CLT2627.sty et de "Séquence 2 - Feuille d'exercices - Vecteurs.tex") :
% sans lui, certaines paires de lettres consécutives dans \cursive
% explosent (une lettre projetée loin à droite), ce qui coupe des mots
% comme "Représenter".
\makeatletter
\ifdefined\XeTeXcharclass
\XeTeXcharclass `a 201 \XeTeXcharclass `b 201 \XeTeXcharclass `c 201
\XeTeXcharclass `d 201 \XeTeXcharclass `e 201 \XeTeXcharclass `f 201
\XeTeXcharclass `g 201 \XeTeXcharclass `h 201 \XeTeXcharclass `i 201
\XeTeXcharclass `j 201 \XeTeXcharclass `k 201 \XeTeXcharclass `l 201
\XeTeXcharclass `m 201 \XeTeXcharclass `n 201 \XeTeXcharclass `o 201
\XeTeXcharclass `p 201 \XeTeXcharclass `q 201 \XeTeXcharclass `r 201
\XeTeXcharclass `s 201 \XeTeXcharclass `t 201 \XeTeXcharclass `u 201
\XeTeXcharclass `v 201 \XeTeXcharclass `w 201 \XeTeXcharclass `x 201
\XeTeXcharclass `y 201 \XeTeXcharclass `z 201
\XeTeXcharclass `A 201 \XeTeXcharclass `B 201 \XeTeXcharclass `C 201
\XeTeXcharclass `D 201 \XeTeXcharclass `E 201 \XeTeXcharclass `F 201
\XeTeXcharclass `G 201 \XeTeXcharclass `H 201 \XeTeXcharclass `I 201
\XeTeXcharclass `J 201 \XeTeXcharclass `K 201 \XeTeXcharclass `L 201
\XeTeXcharclass `M 201 \XeTeXcharclass `N 201 \XeTeXcharclass `O 201
\XeTeXcharclass `P 201 \XeTeXcharclass `Q 201 \XeTeXcharclass `R 201
\XeTeXcharclass `S 201 \XeTeXcharclass `T 201 \XeTeXcharclass `U 201
\XeTeXcharclass `V 201 \XeTeXcharclass `W 201 \XeTeXcharclass `X 201
\XeTeXcharclass `Y 201 \XeTeXcharclass `Z 201
\XeTeXcharclass `\é 201 \XeTeXcharclass `\è 201 \XeTeXcharclass `\ê 201
\XeTeXcharclass `\ë 201 \XeTeXcharclass `\É 201 \XeTeXcharclass `\È 201
\XeTeXcharclass `\Ê 201 \XeTeXcharclass `\Ë 201 \XeTeXcharclass `\à 201
\XeTeXcharclass `\â 201 \XeTeXcharclass `\ä 201 \XeTeXcharclass `\À 201
\XeTeXcharclass `\Â 201 \XeTeXcharclass `\Ä 201 \XeTeXcharclass `\ô 201
\XeTeXcharclass `\ö 201 \XeTeXcharclass `\Ô 201 \XeTeXcharclass `\Ö 201
\XeTeXcharclass `\î 201 \XeTeXcharclass `\ï 201 \XeTeXcharclass `\Î 201
\XeTeXcharclass `\Ï 201 \XeTeXcharclass `\ù 201 \XeTeXcharclass `\û 201
\XeTeXcharclass `\ü 201 \XeTeXcharclass `\Ù 201 \XeTeXcharclass `\Û 201
\XeTeXcharclass `\Ü 201 \XeTeXcharclass `\ç 201 \XeTeXcharclass `\Ç 201
\XeTeXcharclass `\œ 201 \XeTeXcharclass `\Œ 201 \XeTeXcharclass `\æ 201
\XeTeXcharclass `\Æ 201
\let\cursive@orig\cursive
\renewcommand{\cursive}{%
  \cursive@orig
  \XeTeXinterchartoks 201 201 = {\kern0pt }%
}
\fi
\makeatother
\usepackage{tikz}
\usetikzlibrary{calc,arrows.meta,shapes.misc}
\tikzset{ptcross/.style={cross out,draw,thick,minimum size=3pt,inner sep=0pt}}
\usepackage[hidelinks]{hyperref}

\definecolor{diffgreen}{HTML}{40C057}
\definecolor{diffblue}{HTML}{228BE6}
\definecolor{diffred}{HTML}{FA5252}
\definecolor{skyblue}{HTML}{339AF0}
\definecolor{objorange}{HTML}{F76707}
\newcommand{\starfull}[1]{{\color{#1}\ding{72}}}
\newcommand{\starempty}{{\color{gray!45}\ding{73}}}
\newcommand{\Vj}[1]{\overrightarrow{\text{#1}}}
\newcommand{\Pt}[1]{\textsf{\footnotesize #1}}

% ============================================================ %
% Taille de police : 20pt pour le texte courant (un peu moins que les
% 22pt du cours en paysage, la page portrait étant plus étroite), tailles
% voisines redéfinies en proportion.
% ============================================================ %
\renewcommand{\normalsize}{\fontsize{20}{27}\selectfont}
\renewcommand{\small}{\fontsize{17}{22}\selectfont}
\renewcommand{\large}{\fontsize{23}{29}\selectfont}
\renewcommand{\Large}{\fontsize{27}{33}\selectfont}
\renewcommand{\LARGE}{\fontsize{31}{38}\selectfont}
\renewcommand{\huge}{\fontsize{35}{42}\selectfont}
\renewcommand{\Huge}{\fontsize{39}{46}\selectfont}
\normalsize
\linespread{1.15}
\setlength{\parskip}{10pt}
\setlength{\parindent}{0pt}
\RaggedRight % texte non justifié (ragged2e) dans tout le document

\newcommand*{\textcircledd}[1]{\begin{tikzpicture}[baseline=(letter.base)]\node[draw,circle,inner sep=3pt] (letter) {#1};\end{tikzpicture}}
\newcommand{\numeroCerne}[1]{\protect\textcircledd{\bfseries\large #1}}
\setlist[enumerate,1]{label=\numeroCerne{\arabic*}, leftmargin=*, labelsep=10pt, itemsep=9pt, topsep=6pt, partopsep=0pt, parsep=0pt}
\setlist[enumerate,2]{label=(\alph*), leftmargin=1.6em, itemsep=7pt, topsep=5pt, parsep=0pt}
\setlist[itemize,1]{leftmargin=1.4em, itemsep=8pt, topsep=6pt}
% Sous-questions (a), (b), (c)... en ligne plutôt qu'empilées : à cette
% taille de police, une sous-question par ligne fait déborder les
% exercices 11 et 12 de leur page.
\newlist{souslist}{enumerate*}{1}
\setlist[souslist]{label=(\alph*), itemjoin=\hspace{1.6em}, mode=unboxed}

% Titre "Exercice N" + étoiles + compétence cursive, agrandi.
\newcommand{\ExoTitre}[3]{%
\noindent\tcbox[on line, colback=skyblue, colframe=skyblue, boxsep=4pt, left=10pt, right=10pt, top=4pt, bottom=4pt, arc=5pt]{\textcolor{white}{\large\ExoTexte{Exercice~#1}}}\quad{\Large #2}\hfill{\Large\color{skyblue}\cursive\hyphenchar\font=-1 #3}\par
\vspace{16pt}%
}

% Rappel d'objectif en haut de chaque page (répété à chaque exercice).
\newcommand{\ObjRappel}[2]{%
% \sloppy : sans lui, le rappel de l'objectif 1 dépasse dans la marge de
% droite (TeX préfère une ligne trop longue à un renvoi à la ligne).
\begingroup\sloppy
\noindent{\color{objorange}\bfseries\large Objectif #1~:}~{\large\color{objorange}#2}\par
\endgroup
\vspace{6pt}
\noindent{\color{gray!45}\rule{\linewidth}{0.6pt}}\par
\vspace{18pt}%
}

% Mise en page d'un exercice avec figure : chaque exercice tient sur UNE
% page. La page étant en paysage (large mais basse), l'énoncé est placé à
% gauche et la figure à droite, ce qui laisse à la figure toute la
% hauteur restante au lieu de la coincer sous le texte.
%
% La figure reprend le tikzpicture original tel quel (mêmes coordonnées,
% même échelle) et est redimensionnée en bloc, ce qui agrandit
% uniformément le dessin ET les étiquettes \Pt. La hauteur disponible est
% mesurée ici (\pagegoal-\pagetotal = place verticale restante sous le
% titre de l'exercice) plutôt que fixée à la main, puis largeur ET
% hauteur sont contraintes (adjustbox, pas resizebox) : la figure ne peut
% donc jamais basculer seule sur la page suivante.
%
% Le \vspace{0pt} en tête de chaque minipage aligne les deux colonnes par
% leur bord supérieur et non par la ligne de base de leur première ligne ;
% sans lui, la figure s'empilerait au-dessus de cette ligne de base et le
% bloc déborderait de la page.
%
% #1 = largeur de la colonne de texte (fraction de \linewidth, à réduire
% pour les figures larges), #2 = énoncé, #3 = figure.
\newlength{\GCavail}
\newlength{\GCparskip}
% (+m : les arguments doivent accepter les paragraphes, certains textes
%  de corrigé contiennent une ligne blanche.)
\NewDocumentCommand{\GCexo}{O{0.46} +m +m}{%
\par\penalty0\relax
\ifdim\pagegoal=\maxdimen
  \setlength{\GCavail}{\textheight}%
\else
  \setlength{\GCavail}{\dimexpr\pagegoal-\pagetotal\relax}%
\fi
\addtolength{\GCavail}{-0.8\baselineskip}% marge de sécurité en bas de page
\ifdim\GCavail<3cm\relax\setlength{\GCavail}{3cm}\fi
\setlength{\GCparskip}{\parskip}%
\noindent
\begin{minipage}[t]{#1\linewidth}\vspace{0pt}%
% Une minipage réinitialise les paramètres de paragraphe : sans ces deux
% lignes, le texte de la colonne redeviendrait justifié (avec des espaces
% étirés) et les paragraphes ne seraient plus séparés.
\RaggedRight
\setlength{\parskip}{\GCparskip}%
\hyphenpenalty=10000\exhyphenpenalty=10000% pas de mots coupés en fin de ligne

#2
\end{minipage}\hfill
\begin{minipage}[t]{\dimexpr0.96\linewidth-#1\linewidth\relax}\vspace{0pt}%
\centering
\adjustbox{width=\linewidth,height=\GCavail,keepaspectratio}{#3}%
\end{minipage}\par
}

\pagestyle{fancy}
\fancyhf{}
\renewcommand{\headrulewidth}{0pt}
\fancyfoot[C]{\small Feuille d'exercices n\textdegree{}2 --- Vecteurs --- \thepage/\pageref{LastPage}}

\begin{document}

% ============================================================ %
% Page de garde
% ============================================================ %
\thispagestyle{empty}
\begin{center}
\vspace*{2cm}
\begin{tcolorbox}[
  enhanced,
  width=0.95\linewidth,
  colback=white,
  colframe=black,
  boxrule=1pt,
  arc=10pt,
  sharp corners=downhill,
  left=16pt, right=16pt, top=18pt, bottom=12pt,
  halign=center,
  overlay={
    \node[fill=white, anchor=west, font=\large, inner xsep=6pt]
      at ([xshift=18pt]frame.north west) {Feuille d'exercices n\textdegree{}2};
  },
]
{\Huge Vecteurs du plan}
\end{tcolorbox}
\vspace{1.8cm}
{\Large Un exercice par page.}\\[10pt]
{\large Exercices 1 \`a 16 (Objectifs 1 \`a 4).}
\end{center}
\newpage
\setcounter{page}{1}

% ============================================================ %
% Objectif 1
% ============================================================ %

\ObjRappel{1}{Vecteur $\Vj{MM'}$ associ\'e \`a la translation qui transforme $M$ en $M'$ --- direction, sens et norme}
\ExoTitre{1}{\starfull{diffgreen}\starempty\starempty}{Repr\'esenter}
\GCexo{%
\`A partir de la figure ci-contre,
\begin{enumerate}
\item donner les images des points $C$, $D$, $E$ dans la translation de vecteur $\Vj{AB}$~;
\item en d\'eduire trois vecteurs \'egaux au vecteur $\Vj{AB}$.
\end{enumerate}
}{%
\begin{tikzpicture}[scale=0.38]
\draw[thin,gray!35] (0,-1) grid (8,6);
\coordinate (A) at (1,4); \coordinate (B) at (4,5);
\coordinate (C) at (2,2); \coordinate (D) at (4,1); \coordinate (E) at (1,0);
\coordinate (Y) at (5,3); \coordinate (Z) at (7,2); \coordinate (X) at (6,1);
\draw[-{Stealth},thick] (A) -- (B);
\node[left] at (A) {\Pt{A}}; \node[above] at (B) {\Pt{B}};
\foreach \p in {C,D,E,Y,Z,X} \node[ptcross] at (\p) {};
\node[left] at (C) {\Pt{C}}; \node[left] at (D) {\Pt{D}}; \node[above] at (E) {\Pt{E}};
\node[right] at (Y) {\Pt{Y}}; \node[left] at (Z) {\Pt{Z}}; \node[right] at (X) {\Pt{X}};
\end{tikzpicture}}

\newpage
\ExoTitre{2}{\starfull{diffgreen}\starempty\starempty}{Repr\'esenter}
\GCexo[0.40]{%
Sans justifier, donner l'image de chaque point par la translation de vecteur indiqu\'ee.
\begin{enumerate}
\item Image du point $I$ par la translation de vecteur $\Vj{AD}$.
\item Image du point $J$ par la translation de vecteur $\Vj{LR}$.
\item Image du point $R$ par la translation de vecteur $\Vj{NC}$.
\end{enumerate}
}{%
\begin{tikzpicture}[scale=0.32]
\draw[thin,gray!35] (-1,-1) grid (11,5);
\coordinate (A) at (0,4); \coordinate (B) at (2,4); \coordinate (C) at (4,4); \coordinate (D) at (6,4); \coordinate (E) at (8,4); \coordinate (F) at (10,4);
\coordinate (G) at (0,2); \coordinate (H) at (2,2); \coordinate (I) at (4,2); \coordinate (J) at (6,2); \coordinate (K) at (8,2); \coordinate (L) at (10,2);
\coordinate (M) at (0,0); \coordinate (N) at (2,0); \coordinate (O) at (4,0); \coordinate (P) at (6,0); \coordinate (R) at (8,0); \coordinate (S) at (10,0);
\foreach \p in {A,B,C,D,E,F,G,H,I,J,K,L,M,N,O,P,R,S} \node[ptcross] at (\p) {};
\node[above] at (A) {\Pt{A}}; \node[above] at (B) {\Pt{B}}; \node[above] at (C) {\Pt{C}}; \node[above] at (D) {\Pt{D}}; \node[above] at (E) {\Pt{E}}; \node[above] at (F) {\Pt{F}};
\node[above] at (G) {\Pt{G}}; \node[above] at (H) {\Pt{H}}; \node[above] at (I) {\Pt{I}}; \node[above] at (J) {\Pt{J}}; \node[above] at (K) {\Pt{K}}; \node[above] at (L) {\Pt{L}};
\node[above] at (M) {\Pt{M}}; \node[above] at (N) {\Pt{N}}; \node[above] at (O) {\Pt{O}}; \node[above] at (P) {\Pt{P}}; \node[above] at (R) {\Pt{R}}; \node[above] at (S) {\Pt{S}};
\end{tikzpicture}}

\newpage
\ExoTitre{3}{\starfull{diffblue}\starfull{diffblue}\starempty}{Repr\'esenter}
\GCexo{%
La figure ci-contre repr\'esente sept hexagones r\'eguliers et num\'erot\'es.
Compl\`ete les phrases suivantes~:
\begin{enumerate}
\item L'hexagone \dots{} est l'image de l'hexagone 1 dans la translation de vecteur $\Vj{AC}$.
\item L'hexagone 6 est l'image de l'hexagone \dots{} dans la translation de vecteur $\Vj{AB}$.
\item L'hexagone 3 est l'image de l'hexagone \dots{} dans la translation de vecteur $\Vj{DE}$.
\end{enumerate}
}{%
\begin{tikzpicture}[scale=0.22]
\draw (2,0)-- (4,0)-- (5,1.73)-- (4,3.46)-- (2,3.46)-- (1,1.73)-- cycle;
\draw (3,1.73) node[draw,circle,inner sep=1pt]{\footnotesize 1};
\draw (4,3.46)-- (5,5.2)-- (4,6.93)-- (2,6.93)-- (1,5.2)-- (2,3.46);
\draw (3,5.2) node[draw,circle,inner sep=1pt]{\footnotesize 2};
\draw (5,1.73) --(7,1.73)-- (8,3.46)-- (7,5.2)-- (5,5.2);
\draw (6,3.46) node[draw,circle,inner sep=1pt]{\footnotesize 3};
\draw (4,0)-- (5,-1.73)-- (7,-1.73)-- (8,0)--(7,1.73);
\draw (6,0) node[draw,circle,inner sep=1pt]{\footnotesize 4};
\draw (8,0)-- (10,0)-- (11,1.73)-- (10,3.46)-- (8,3.46);
\draw (9,1.73) node[draw,circle,inner sep=1pt]{\footnotesize 5};
\draw (10,3.46)-- (11,5.2)-- (10,6.93)-- (8,6.93)-- (7,5.2);
\draw (9,5.2) node[draw,circle,inner sep=1pt]{\footnotesize 6};
\draw (8,6.93)-- (7,8.66)-- (5,8.66)-- (4,6.93);
\draw (6,6.93) node[draw,circle,inner sep=1pt]{\footnotesize 7};
\coordinate (A) at (2,3.46); \coordinate (B) at (5,8.66); \coordinate (C) at (8,6.93); \coordinate (D) at (11,5.2); \coordinate (E) at (11,1.73);
\foreach \p in {A,B,C,D,E} \node[ptcross] at (\p) {};
\node[left] at (A) {\Pt{A}};
\node[left] at (B) {\Pt{B}};
\node[above right] at (C) {\Pt{C}};
\node[right] at (D) {\Pt{D}};
\node[right] at (E) {\Pt{E}};
\end{tikzpicture}}

\newpage
\ExoTitre{4}{\starfull{diffred}\starfull{diffred}\starfull{diffred}}{Raisonner}
Soit $\vec u$ un vecteur non nul, et $A$ un point. On note $A'$ l'image de $A$ par la translation de vecteur $\vec u$.
\begin{enumerate}
\item On consid\`ere un point $B$ tel que $\Vj{AB}$ ait la m\^eme direction et le m\^eme sens que $\vec u$, mais une norme diff\'erente de celle de $\vec u$. Le point $B$ peut-il \^etre \'egal \`a $A'$~? Justifier.
\item M\^eme question si $\Vj{AB}$ a la m\^eme direction que $\vec u$ mais le sens contraire.
\item \'Enoncer, \`a l'aide des trois caract\'eristiques d'un vecteur, une condition n\'ecessaire et suffisante pour que $B=A'$.
\end{enumerate}

% ============================================================ %
% Objectif 2
% ============================================================ %

\newpage
\ObjRappel{2}{\'Egalit\'e de deux vecteurs, vecteur nul}
\ExoTitre{5}{\starfull{diffgreen}\starempty\starempty}{Raisonner}
\GCexo{%
Sur la figure ci-contre, expliquer, en utilisant les termes direction, sens ou norme, pourquoi le vecteur $\Vj{AB}$ n'est \'egal \`a aucun des autres vecteurs repr\'esent\'es.
}{%
\begin{tikzpicture}[scale=0.4]
\draw[-{Stealth},thick] (0,2) -- (1,3); \node[below] at (0,2) {\Pt{A}}; \node[above] at (1,3) {\Pt{B}};
\draw[-{Stealth},thick] (1,0) -- (3,0); \node[below] at (1,0) {\Pt{C}}; \node[below] at (3,0) {\Pt{D}};
\draw[-{Stealth},thick] (3,2) -- (5,4); \node[below] at (3,2) {\Pt{G}}; \node[above] at (5,4) {\Pt{H}};
\draw[-{Stealth},thick] (7,2) -- (6,1); \node[below] at (6,1) {\Pt{E}}; \node[above] at (7,2) {\Pt{F}};
\end{tikzpicture}}

\newpage
\ExoTitre{6}{\starfull{diffgreen}\starempty\starempty}{Repr\'esenter}
\GCexo{%
\`A partir de la figure ci-contre, citer un vecteur~:
\begin{itemize}[label=$\star$]
\item opposé \`a $\Vj{CD}$
\item m\^eme direction et sens que $\Vj{AC}$
\item m\^eme direction que $\Vj{BC}$ mais sens contraire~
\item \'egal au vecteur $\Vj{BA}$.
\end{itemize}
}{%
\begin{tikzpicture}[scale=0.32]
\draw[thin,gray!35] (-0.5,-0.5) grid (9.5,8.5);
\coordinate (A) at (3,4); \coordinate (B) at (1,1); \coordinate (C) at (3,1); \coordinate (D) at (4,2);
\foreach \p in {A,B,C,D} \node[ptcross] at (\p) {};
\node[left] at (A) {\Pt{A}}; \node[below] at (B) {\Pt{B}}; \node[below] at (C) {\Pt{C}}; \node[below] at (D) {\Pt{D}};
\draw[-{Stealth},thick] (0,3) -- (1,3) node[midway,above]{\footnotesize $\vec u$};
\draw[-{Stealth},thick] (1,5) -- (3,8) node[midway,above left]{\footnotesize $\vec v$};
\draw[-{Stealth},thick] (4,5) -- (4,7) node[midway,left]{\footnotesize $\vec w$};
\draw[-{Stealth},thick] (5,6) -- (6,7) node[midway,above]{\footnotesize $\vec r$};
\draw[-{Stealth},thick] (8,8) -- (6,3) node[midway,above]{\footnotesize $\vec s$};
\draw[-{Stealth},thick] (8,8) -- (7,8) node[midway,above]{\footnotesize $\vec t$};
\draw[-{Stealth},thick] (9,2) -- (8,1) node[midway,above]{\footnotesize $\vec p$};
\draw[-{Stealth},thick] (9,5) -- (9,4) node[midway,left]{\footnotesize $\vec m$};
\end{tikzpicture}}

\newpage
\ExoTitre{7}{\starfull{diffblue}\starfull{diffblue}\starempty}{Raisonner}
Soient $A$, $B$, $C$ et $D$ quatre points du plan tels que $\Vj{AB}=\Vj{CD}$.

D\'emontrer que $\Vj{AC}=\Vj{BD}$.

\newpage
\ExoTitre{8}{\starfull{diffred}\starfull{diffred}\starfull{diffred}}{Raisonner}
Soient $ABCD$ et $ABEF$ deux parall\'elogrammes.
\begin{enumerate}
\item R\'ealiser une figure \`a main lev\'ee.
\item Donner deux vecteurs \'egaux au vecteur $\Vj{AB}$.
\item En d\'eduire la nature du quadrilat\`ere $EFDC$.
\end{enumerate}

\newpage
\ExoTitre{9}{\starfull{diffred}\starfull{diffred}\starfull{diffred}}{Raisonner}
\GCexo{%
Sur la figure ci-contre, $ABCD$ est un rectangle, et $AEBC$ et $ACDF$ sont des parall\'elogrammes.
\begin{enumerate}
\item Montrer que $\Vj{EB}=\Vj{FD}$.
\item Montrer que le quadrilat\`ere $EBDF$ est un parall\'elogramme de centre $A$.
\end{enumerate}
}{%
\begin{tikzpicture}[scale=0.5]
\coordinate (D) at (0,2); \coordinate (C) at (3,2); \coordinate (A) at (0,0); \coordinate (B) at (3,0);
\coordinate (F) at (-3,0); \coordinate (E) at (0,-2);
\draw[blue] (D)--(C)--(B)--(A)--cycle;
\draw[green!50!black] (F)--(A)--(C)--(D)--cycle;
\draw[red] (F)--(A)--(E)--(B);
\node[above] at (D) {\Pt{D}}; \node[above] at (C) {\Pt{C}}; \node[below left] at (A) {\Pt{A}}; \node[below] at (B) {\Pt{B}};
\node[left] at (F) {\Pt{F}}; \node[below] at (E) {\Pt{E}};
\end{tikzpicture}}

% ============================================================ %
% Objectif 3
% ============================================================ %

\newpage
\ObjRappel{3}{Repr\'esentants d'un vecteur}
\ExoTitre{10}{\starfull{diffgreen}\starempty\starempty}{Repr\'esenter}
\GCexo{%
Sur la figure ci-contre,
\begin{enumerate}
\item citer deux repr\'esentants du vecteur $\Vj{AB}$, autres que lui-m\^eme~;
\item indiquer si le vecteur $\Vj{EF}$ est aussi un repr\'esentant du vecteur $\Vj{AB}$, en justifiant.
\end{enumerate}
}{%
\begin{tikzpicture}[scale=0.42]
\draw[thin,gray!35] (-0.5,-0.5) grid (7.5,6.5);
\coordinate (A) at (0,1); \coordinate (B) at (2,2);
\coordinate (C) at (3,4); \coordinate (D) at (5,5);
\coordinate (G) at (5,0); \coordinate (H) at (7,1);
\coordinate (E) at (1,4); \coordinate (F) at (2,6);
\draw[-{Stealth},thick,blue] (A) -- (B); \draw[-{Stealth},thick,blue] (C) -- (D); \draw[-{Stealth},thick,blue] (G) -- (H);
\draw[-{Stealth},thick,red] (E) -- (F);
\node[below] at (A) {\Pt{A}}; \node[right] at (B) {\Pt{B}};
\node[left] at (C) {\Pt{C}}; \node[right] at (D) {\Pt{D}};
\node[below] at (G) {\Pt{G}}; \node[right] at (H) {\Pt{H}};
\node[left] at (E) {\Pt{E}}; \node[above] at (F) {\Pt{F}};
\end{tikzpicture}}

\newpage
\ExoTitre{11}{\starfull{diffgreen}\starempty\starempty}{Calculer}
$EFGH$ est un parall\'elogramme de centre $I$.
\begin{enumerate}
\item Donner un vecteur \'egal \`a~:
\begin{souslist}
\item $\Vj{EF}$
\item $\Vj{GH}$
\item $\Vj{EH}$
\item $\Vj{GF}$
\end{souslist}
\item Donner un vecteur oppos\'e au vecteur $\Vj{EF}$.
\item Compl\'eter par le point qui convient~:
\begin{souslist}
\item $\Vj{EI}=\Vj{I\dots}$
\item $\Vj{HI}=\Vj{I\dots}$
\item $\Vj{FI}=\Vj{I\dots}$
\end{souslist}
\item Compl\'eter~: $\Vj{EH}=-\Vj{\vphantom{H}\dots}$
\item Quel est le repr\'esentant d'origine $F$~:
\begin{souslist}
\item du vecteur $\Vj{EH}$~?
\item du vecteur $\Vj{IH}$~?
\end{souslist}
\end{enumerate}

\newpage
\ExoTitre{12}{\starfull{diffblue}\starfull{diffblue}\starempty}{Repr\'esenter}
\GCexo[0.53]{%
On utilisera la figure ci-contre~:
{\setlist[enumerate,1]{label=\numeroCerne{\arabic*}, leftmargin=*, labelsep=10pt, itemsep=3pt, topsep=2pt, partopsep=0pt, parsep=0pt}
\setlist[enumerate,2]{label=(\alph*), leftmargin=1.6em, itemsep=1pt, topsep=1pt, parsep=0pt}
\begin{enumerate}
\item Citer le repr\'esentant du vecteur $\vec u$~:
\begin{souslist}
\item d'origine $A$
\item d'origine $C$
\end{souslist}
\item Citer le repr\'esentant du vecteur $\vec u$~:
\begin{souslist}
\item d'extr\'emit\'e $D$
\item d'extr\'emit\'e $F$
\end{souslist}
\item Citer le repr\'esentant du vecteur $\vec v$~:
\begin{souslist}
\item d'origine $P$
\item d'origine $S$
\end{souslist}
\item Citer le repr\'esentant du vecteur $\vec v$~:
\begin{souslist}
\item d'extr\'emit\'e $H$
\item d'extr\'emit\'e $W$
\end{souslist}
\end{enumerate}}
}{%
\begin{tikzpicture}[scale=0.42]
\draw[thin,gray!35] (-0.5,-0.5) grid (7.5,7.5);
\draw[-{Stealth},thick,blue] (4,6) -- (6,7) node[midway,above]{\footnotesize $\vec u$};
\draw[-{Stealth},thick,teal] (2,7) -- (4,5) node[midway,above]{\footnotesize $\vec v$};
\coordinate (A) at (1,0); \coordinate (K) at (3,1);
\coordinate (C) at (5,3); \coordinate (D) at (7,4);
\coordinate (M) at (0,4); \coordinate (F) at (2,5);
\coordinate (G) at (4,0); \coordinate (J) at (0,1);
\coordinate (N) at (6,2); \coordinate (R) at (1,3);
\coordinate (T) at (4,4);
\coordinate (P) at (2,7); \coordinate (Q) at (4,5);
\coordinate (S) at (4,7); \coordinate (H) at (6,5);
\coordinate (U) at (4,3); \coordinate (W) at (6,1);
\foreach \p in {A,K,C,D,M,F,G,J,N,R,T,P,Q,S,H,U,W} \node[ptcross] at (\p) {};
\node[below] at (A) {\Pt{A}}; \node[above] at (K) {\Pt{K}};
\node[right] at (C) {\Pt{C}}; \node[above] at (D) {\Pt{D}};
\node[left] at (M) {\Pt{M}}; \node[above] at (F) {\Pt{F}};
\node[below] at (G) {\Pt{G}}; \node[left] at (J) {\Pt{J}};
\node[right] at (N) {\Pt{N}}; \node[left] at (R) {\Pt{R}};
\node[left] at (T) {\Pt{T}};
\node[above] at (P) {\Pt{P}}; \node[right] at (Q) {\Pt{Q}};
\node[above] at (S) {\Pt{S}}; \node[below] at (H) {\Pt{H}};
\node[below] at (U) {\Pt{U}}; \node[below] at (W) {\Pt{W}};
\end{tikzpicture}}

% ============================================================ %
% Objectif 4
% ============================================================ %

\newpage
\ObjRappel{4}{Somme de deux vecteurs et relation de Chasles}
\ExoTitre{13}{\starfull{diffgreen}\starempty\starempty}{Repr\'esenter}
\GCexo[0.44]{%
Construire~:
\begin{enumerate}
\item le point $P$ tel que $\Vj{BP}=\vec u+\vec v$~;
\item le point $G$ tel que $\Vj{VG}=\vec v+\vec w$~;
\item le point $K$ tel que $\Vj{IK}=\vec w+\vec u$.
\end{enumerate}
}{%
\begin{tikzpicture}[scale=0.5]
\draw[thin,black!70] (0,0) grid (12,6);
\draw[-{Stealth},thick,blue] (2,5) -- (2,3) node[midway,right]{\footnotesize $\vec u$};
\draw[-{Stealth},thick,red] (6,5) -- (8,4) node[midway,above]{\footnotesize $\vec v$};
\draw[-{Stealth},thick,teal] (4,1) -- (6,2) node[midway,below]{\footnotesize $\vec w$};
\coordinate (B) at (1,4); \node[ptcross] at (B) {}; \node[left] at (B) {\Pt{B}};
\coordinate (V) at (4,4); \node[ptcross] at (V) {}; \node[below] at (V) {\Pt{V}};
\coordinate (I) at (1,1); \node[ptcross] at (I) {}; \node[below] at (I) {\Pt{I}};
\end{tikzpicture}}

\newpage
\ExoTitre{14}{\starfull{diffgreen}\starempty\starempty}{Calculer}
\'Ecrire chacune des sommes suivantes \`a l'aide d'un seul vecteur, en utilisant la relation de Chasles.
\begin{enumerate}
\item $\Vj{TR}+\Vj{UT}+\Vj{RU}$
\item $\Vj{ST}-\Vj{VU}+\Vj{TU}$
\end{enumerate}

\newpage
\ExoTitre{15}{\starfull{diffblue}\starfull{diffblue}\starempty}{Repr\'esenter}
\GCexo{%
$ABCDEF$ est un hexagone r\'egulier de centre $O$.
Compl\'eter les \'egalit\'es suivantes en n'utilisant que des noms de points de la figure.
\begin{enumerate}
\item $\Vj{AB}+\Vj{CD}=\dots$
\item $\Vj{OB}+\Vj{OE}=\dots$
\item $\Vj{FO}+\Vj{DO}=\dots$
\item $\Vj{CA}-\Vj{DC}=\dots$
\item $\Vj{AF}+\Vj{CD}=\dots$
\item $\Vj{AB}-\Vj{OA}+\Vj{AF}=\dots$
\end{enumerate}
}{%
\begin{tikzpicture}[scale=0.55]
\foreach \i/\lab/\pos in {0/A/below right,60/B/above right,120/C/above,180/D/above left,240/E/below left,300/F/below}
\coordinate (\lab) at (\i:1.4);
\draw (A)--(B)--(C)--(D)--(E)--(F)--cycle;
\draw (A)--(D); \draw (B)--(E); \draw (C)--(F);
\node[below right] at (A) {\Pt{A}}; \node[above right] at (B) {\Pt{B}}; \node[above] at (C) {\Pt{C}};
\node[above left] at (D) {\Pt{D}}; \node[below left] at (E) {\Pt{E}}; \node[below] at (F) {\Pt{F}};
\node[anchor=north] at (0,0) {\Pt{O}};
\end{tikzpicture}}

\newpage
\ExoTitre{16}{\starfull{diffblue}\starfull{diffblue}\starempty}{Raisonner}
Soit $ABC$ un triangle et $I$ le milieu de $[BC]$.
\begin{enumerate}
\item Justifier que $\Vj{IB}+\Vj{IC}=\vec 0$.
\item En utilisant la relation de Chasles, d\'emontrer que~: $\Vj{AB}+\Vj{AC}=2\Vj{AI}$.
\end{enumerate}

\end{document}
